\documentclass[11pt]{book}
\usepackage[margin=1in]{geometry}
\usepackage{amsmath,amssymb}
\usepackage{graphicx}
\usepackage{tikz-cd}
\usepackage{multicol}

\setlength{\parindent}{0pt}
\setlength{\parskip}{0.75\baselineskip}

\begin{document}

\chapter{Low-Rank Structure in Transformer Gradient Matrices}

\section{Motivation and Overview}

Modern transformer models contain large linear layers whose parameter matrices
(e.g., $W_q$, $W_k$, $W_v$, and $W_o$) often have dimensions on the order of thousands.
During training, the gradients of these matrices exhibit strong regularities:
they are neither random nor isotropic.
Empirically, a small number of dominant directions account for a large fraction
of the gradient energy, while the remaining directions are highly diverse and
context-dependent.

Understanding and exploiting this structure is important for several reasons.
First, it sheds light on why optimization dynamics in large models are often
dominated by a few global factors.
Second, it motivates algorithmic simplifications, such as low-rank or
subspace-based optimizers, which can reduce memory and computation.
Third, it provides a principled explanation for why very simple operations,
such as averaging, can sometimes extract meaningful information from high-dimensional
gradients.

This chapter develops a simple but rigorous framework for understanding these
phenomena. We focus on gradients of transformer linear layers and show how
averaging operations can be interpreted as probes that isolate shared,
grammar-like structure while suppressing diverse semantic effects.
We then generalize this idea from rank-one approximations to low-rank
($k$-dimensional) subspaces.

\section{Gradients of Linear Layers in Transformers}

Consider a single transformer layer and one of its linear projections, for example
the key projection matrix
\[
W_k \in \mathbb{R}^{d \times d}.
\]
Given an input hidden state $x \in \mathbb{R}^d$, the key vector is computed as
$k = W_k x$.

During backpropagation, the gradient of the loss $\mathcal{L}$ with respect to $W_k$
has the same shape as $W_k$:
\[
G \;=\; \nabla_{W_k} \mathcal{L} \in \mathbb{R}^{d \times d}.
\]
The entry $G_{i\ell}$ should be read concretely as follows: it measures how the loss
would change if we increase the weight that connects the $\ell$-th input feature
to the $i$-th output feature of the key projection.

For a single token (or training example) indexed by $j$, the contribution to the
gradient has an outer-product form:
\[
\nabla_{W_k} \mathcal{L}^{(j)} = \delta_k^{(j)} (x^{(j)})^\top,
\]
where $x^{(j)} \in \mathbb{R}^d$ is the input activation for that token and
$\delta_k^{(j)} \in \mathbb{R}^d$ is the backpropagated error signal at the key
projection. The full gradient over $n$ tokens is therefore
\[
G = \sum_{j=1}^{n} \delta_k^{(j)} (x^{(j)})^\top.
\]
Thus, although $G$ is a single $d \times d$ matrix, it implicitly aggregates
information from many tokens and contexts.

\section{Common Structure vs.\ Semantic Variability}

Not all token-level gradient contributions are alike.
Some correspond to broadly shared linguistic structure, such as syntax,
function words, or positional patterns.
These effects tend to appear across many tokens and therefore push the parameters
in similar directions.
Other contributions reflect specific meanings, rare words, or unusual contexts.
These semantic effects vary widely from token to token.

When summed together, shared structural contributions reinforce one another,
while diverse semantic contributions often interfere destructively.
This motivates a conceptual decomposition
\[
G = A + R,
\]
where $A$ represents the \emph{common} (grammar-like) component, and $R$
represents the \emph{semantic/long-tail} remainder.

\section{Problem Formulation}

We are given $G \in \mathbb{R}^{d \times d}$ and seek a low-rank approximation that
captures the common structure $A$.
Unlike classical low-rank approximation, our goal is not to faithfully reconstruct
every entry of $G$. Instead, we want to identify directions that are repeatedly
reinforced across many tokens---the directions that dominate optimization dynamics.

\section{Averaging as a Structural Probe}

The simplest linear operation you can perform on a matrix is to sum along one axis.
When you sum the \emph{columns} of $G$, you collapse all input-feature contributions
into a single vector in the output-feature space. When you sum the \emph{rows} of $G$,
you collapse all output-feature contributions into a single vector in the input-feature space.
These two summaries will be the building blocks of a rank-one approximation.

\subsection{Column averaging: what does $u = G\mathbf{1}$ mean?}

Let $\mathbf{1}\in\mathbb{R}^d$ be the all-ones vector. Define
\[
u \stackrel{\mathrm{def}}{=} G\mathbf{1} \in \mathbb{R}^d.
\]
By expanding the matrix--vector product, we get an index-level definition:
\[
u_i = \sum_{\ell=1}^d G_{i\ell}.
\]
So $u_i$ is obtained by fixing an output coordinate $i$ (a row of $G$) and adding
the gradients over \emph{all} input coordinates $\ell$.

Interpretation: $u$ tells us which output directions of $W_k$ are being pushed
consistently \emph{after we forget which input coordinate caused the push}.
In other words, $u$ is a crude ``how much does each output feature want to change?''
summary.

Using the token outer-product form $G=\sum_{j=1}^n \delta_k^{(j)}(x^{(j)})^\top$,
\[
u = G\mathbf{1}
= \sum_{j=1}^n \delta_k^{(j)} \bigl( (x^{(j)})^\top \mathbf{1} \bigr).
\]
Thus $u$ is a weighted sum of the error vectors $\delta_k^{(j)}$, where the weight
\(
(x^{(j)})^\top \mathbf{1}
\)
is simply the sum of the components of the token's input vector.

\subsection{Row averaging: what does $v = G^\top\mathbf{1}$ mean?}

Define
\[
v \stackrel{\mathrm{def}}{=} G^\top \mathbf{1} \in \mathbb{R}^d.
\]
Again, expanding gives
\[
v_\ell = \sum_{i=1}^d G_{i\ell}.
\]
So $v_\ell$ is obtained by fixing an input coordinate $\ell$ (a column of $G$)
and adding the gradients over \emph{all} output coordinates $i$.

Interpretation: $v$ tells us which input-feature directions tend to participate
in strong updates \emph{after we forget which output coordinate they affect}.
In other words, $v$ is a crude ``which input features are consistently involved?''
summary.

In token form,
\[
v = G^\top\mathbf{1}
= \sum_{j=1}^n x^{(j)} \bigl( (\delta_k^{(j)})^\top \mathbf{1} \bigr),
\]
so $v$ is a weighted sum of input activations $x^{(j)}$, weighted by the total
error mass $(\delta_k^{(j)})^\top\mathbf{1}$.

\subsection{Rank-one approximation from row/column averages}

A rank-one matrix is an outer product of two vectors. Since $u$ summarizes the
output-side direction and $v$ summarizes the input-side direction, a natural
rank-one approximation is
\[
\widehat{G}_1 \stackrel{\mathrm{def}}{=} \frac{u v^\top}{\mathbf{1}^\top G \mathbf{1}},
\qquad \text{where } u=G\mathbf{1},\ v=G^\top\mathbf{1}.
\]
The denominator $\mathbf{1}^\top G \mathbf{1}$ is the sum of all entries of $G$.
This scaling ensures that $\widehat{G}_1$ matches that same total sum:
\(
\mathbf{1}^\top \widehat{G}_1 \mathbf{1} = \mathbf{1}^\top G \mathbf{1}.
\)

\section{Assumptions Underlying the Averaging Approximation}

We now state assumptions that make the averaging approximation correct in a precise sense.

Assume the gradient decomposes as
\[
G = A + R,
\]
where $A$ is the shared (grammar-like) component and $R$ is the residual
(semantic/long-tail) component.

\paragraph{Rank-one common structure.}
Assume $A$ is rank one:
\[
A = \sigma a b^\top
\]
for some nonzero vectors $a,b\in\mathbb{R}^d$ and scalar $\sigma>0$.

\paragraph{Coherence with averaging.}
Define the scalars
\[
\alpha \stackrel{\mathrm{def}}{=} \mathbf{1}^\top a,
\qquad
\beta \stackrel{\mathrm{def}}{=} \mathbf{1}^\top b.
\]
Assume $\alpha\neq 0$ and $\beta\neq 0$.
This prevents the common structure from cancelling out under the particular probe
direction $\mathbf{1}$.

\paragraph{Cancellation of the residual under averaging.}
Assume that row- and column-sums of $R$ are small relative to those of $A$:
\[
\|R\mathbf{1}\| \le \varepsilon \|A\mathbf{1}\|,
\qquad
\|R^\top\mathbf{1}\| \le \varepsilon \|A^\top\mathbf{1}\|,
\qquad
|\mathbf{1}^\top R\mathbf{1}| \le \varepsilon |\mathbf{1}^\top A\mathbf{1}|,
\]
for some $0\le \varepsilon<1$.

\section{Theoretical Analysis: Rank-One Case (Complete Proof)}

\paragraph{Theorem.}
Let $u=G\mathbf{1}$, $v=G^\top\mathbf{1}$, $s=\mathbf{1}^\top G \mathbf{1}$, and
\[
\widehat{G}_1 = \frac{u v^\top}{s}.
\]
Under the assumptions above, define
\[
c \stackrel{\mathrm{def}}{=} \frac{\sigma^2 \alpha \beta}{s}.
\]
Then
\[
\|\widehat{G}_1 - c\,ab^\top\|_2
\;\le\;
\frac{2\varepsilon+\varepsilon^2}{1-\varepsilon}\,\|A\|_2,
\qquad
\text{and}\qquad
\left|\frac{c}{\sigma}-1\right|\le \frac{\varepsilon}{1-\varepsilon}.
\]

\paragraph{Proof.}
Because $G=A+R$ and $A=\sigma ab^\top$,
\[
u = G\mathbf{1} = A\mathbf{1} + R\mathbf{1}
= \sigma a(b^\top\mathbf{1}) + r_u
= \sigma\beta\, a + r_u,
\]
where $r_u \stackrel{\mathrm{def}}{=} R\mathbf{1}$.
Similarly,
\[
v = G^\top\mathbf{1} = A^\top\mathbf{1} + R^\top\mathbf{1}
= \sigma b(a^\top\mathbf{1}) + r_v
= \sigma\alpha\, b + r_v,
\]
where $r_v \stackrel{\mathrm{def}}{=} R^\top\mathbf{1}$.
Also,
\[
s=\mathbf{1}^\top G\mathbf{1} = \mathbf{1}^\top A\mathbf{1} + \mathbf{1}^\top R\mathbf{1}
= \sigma\alpha\beta + r_s,
\]
where $r_s\stackrel{\mathrm{def}}{=}\mathbf{1}^\top R\mathbf{1}$.
The bound $|r_s|\le \varepsilon|\sigma\alpha\beta|$ implies
$|s|\ge (1-\varepsilon)|\sigma\alpha\beta|$.

Now expand
\[
\widehat{G}_1 = \frac{uv^\top}{s}
= \frac{(\sigma\beta a + r_u)(\sigma\alpha b + r_v)^\top}{s}.
\]
Subtracting $c\,ab^\top$ and expanding gives
\[
\widehat{G}_1 - c\,ab^\top
= \frac{\sigma\beta\, a r_v^\top}{s}
+ \frac{\sigma\alpha\, r_u b^\top}{s}
+ \frac{r_u r_v^\top}{s}.
\]
Using $\|xy^\top\|_2=\|x\|_2\|y\|_2$ and the assumptions
$\|r_u\|\le \varepsilon\sigma|\beta|\|a\|$ and
$\|r_v\|\le \varepsilon\sigma|\alpha|\|b\|$,
together with $|s|\ge (1-\varepsilon)|\sigma\alpha\beta|$, we obtain
\[
\|\widehat{G}_1 - c\,ab^\top\|_2
\le
\frac{\varepsilon}{1-\varepsilon}\sigma\|a\|\|b\|
+\frac{\varepsilon}{1-\varepsilon}\sigma\|a\|\|b\|
+\frac{\varepsilon^2}{1-\varepsilon}\sigma\|a\|\|b\|.
\]
Since $\|A\|_2=\|\sigma ab^\top\|_2=\sigma\|a\|\|b\|$, the first inequality follows.

Finally,
\[
\frac{c}{\sigma}=\frac{\sigma\alpha\beta}{\sigma\alpha\beta+r_s}
=\frac{1}{1+r_s/(\sigma\alpha\beta)},
\]
and $|r_s/(\sigma\alpha\beta)|\le\varepsilon$ yields
$\left|\frac{c}{\sigma}-1\right|\le \frac{\varepsilon}{1-\varepsilon}$.
\hfill$\square$

\section{Beyond Rank One: Low-Dimensional Grammar Subspaces}

In practice, the common component may have rank $k>1$ (for example, $k\approx 0.01d$).
A single probe direction $\mathbf{1}$ can only produce one mixture direction.
To recover a $k$-dimensional common subspace, we use $k$ independent probe directions.

\section{Rank-$k$ Averaging-Based Approximation (Range Finder)}

Let $\Omega \in \mathbb{R}^{d \times k}$ be a random probe matrix (e.g., Gaussian
or random $\pm 1$ entries). Define the sketch
\[
Y = G \Omega \in \mathbb{R}^{d \times k},
\]
and let $Q\in\mathbb{R}^{d\times k}$ be an orthonormal basis for the columns of $Y$
computed via thin QR. The standard rank-$k$ projected approximation is
\[
\widehat{G}_k = Q Q^\top G.
\]
This is the randomized range-finder method. It generalizes the \emph{idea} of averaging
by using many probe directions, but unlike the rank-one estimator, computing $QQ^\top G$
requires access to $G$ again (a second pass) unless one uses a one-pass sketch
reconstruction.

\section{Summary and Outlook}

Column averaging $u=G\mathbf{1}$ and row averaging $v=G^\top\mathbf{1}$ are simple,
interpretable probes that collapse a gradient matrix along one axis and reveal
directions that are repeatedly reinforced across tokens.
Under a rank-one common-structure model with cancellation of the residual under
averaging, the outer product of these averages yields a principled rank-one
approximation with an explicit $O(\varepsilon)$ error bound.
For higher-rank common structure, multiple probe directions lead naturally to
randomized subspace methods.
\end{document}